
The most-studied benchmarks in machine learning may systematically induce worse generalization. This work demonstrates that models trained on high-popularity datasets exhibit generalization gaps 21 percentage points larger than those trained on less-popular alternatives from the same domain. A ResNet-18 achieving 87.92\% accuracy on CIFAR-10 drops to 69.06\% on CINIC-10, while an identical architecture trained on SVHN achieves 95.32\% and improves to 97.68\% on SVHN-Extra. If the benchmarks relied upon to measure progress systematically induce worse generalization, the field's primary progress indicators may be unreliable.

The surface-level understanding of this phenomenon points to domain shift and distribution mismatch as causes of performance degradation. However, a deeper examination reveals a more concerning pattern: popular benchmarks attract disproportionate research investment. The present analysis finds 7.56 times more optimization-focused papers for high-use datasets compared to low-use alternatives (p=0.027). This creates what is termed the \textit{benchmark co-evolution effect}, where the ML ecosystem's architecture search, hyperparameter tuning, and model design may converge on dataset-specific patterns rather than generalizable features.

The gap in existing work is notable: while Recht et al. (2019) documented 10-15\% accuracy drops on ImageNet with new test sets, and D'Amour et al. (2020) provided theoretical grounding via underspecification, no systematic study has linked cross-repository dataset popularity to generalization failure. This work addresses this gap by quantifying popularity via arXiv paper counts, measuring generalization gaps across canonical dataset pairs, and testing a proposed co-evolution mechanism.

The key empirical finding is that high-popularity datasets (CIFAR-10) exhibit generalization gaps to held-out same-domain datasets that are 21.21 percentage points larger than low-popularity datasets (SVHN), despite identical training protocols and architectures. The first step of the co-evolution mechanism is verified---popular benchmarks attract significantly more optimization research (7.56x, p=0.027)---while texture bias is ruled out as the causal pathway, with modern architectures showing \textit{lower} texture bias than legacy ones.

This work makes the following contributions:

\begin{enumerate}
\item \textbf{Empirical demonstration of the popularity-gap correlation}: High-popularity datasets correlate with larger generalization failures across two canonical dataset pairs, with an effect size (21.21 pp) that is practically significant.
\end{enumerate}

\begin{enumerate}
\item \textbf{Quantification of research investment differential}: A cross-repository analysis linking dataset usage to arXiv optimization paper counts finds a 7.56:1 ratio between high-use and low-use datasets.
\end{enumerate}

\begin{enumerate}
\item \textbf{Mechanism falsification}: The texture bias hypothesis is tested and refuted at CIFAR scale, finding that skip connections in modern architectures improve rather than degrade shape-based generalization.
\end{enumerate}

